% Generated by scripts/assemble.py; corrections belong in source/corrections.json.
\examhead{180}{西安交通大学 2025 年数学分析试卷}{169}
% p170-b24; source PDF page 170
\par\Needspace{5\baselineskip}
\blockanchor{p170-b24}
\majorquestion{一}填空题. 每题 6 分, 共 60 分.\par
% p170-b25; source PDF page 170
\blockanchor{p170-b25}
\question{1}设  $x_{n}>0,\lim_{n\to\infty}x_{n}=1$ ，则  $\lim_{n\to\infty}\sqrt[n]{x_{1}x_{2}\cdots x_{n}}=$  \_\_\_\_.\par
% p170-b26; source PDF page 170
\blockanchor{p170-b26}
\question{2}设  $f(x)=x^{2}e^{x}$ ，则  $f^{(10)}(0)=$  \_\_\_\_.\par
% p171-b00; source PDF page 171
\blockanchor{p171-b00}
\question{3}$\lim_{x\to0}\left(\frac{2^{x^{2}}+1}{2^{x}+1}\right)^{\frac{1}{x}}=$  \_\_\_\_.\par
% p171-b01; source PDF page 171
\blockanchor{p171-b01}
\question{4}$\int_{0}^{\frac{\pi}{2}}\ln \sin x\mathrm{d}x =$\par
% p171-b02; source PDF page 171
\blockanchor{p171-b02}
\question{5}函数 $f(x, y) = x^3 - x^2 - xy + y^2$ 的极值点为\par
% p171-b03; source PDF page 171
\par\Needspace{6\baselineskip}
\blockanchor{p171-b03}
\question{6}已知\par
\nopagebreak[4]
% p171-b04; source PDF page 171
\blockanchor{p171-b04}
\displayformula{f (x, y) = \left\{ \begin{array}{l l} x y \frac {x ^ {2} - y ^ {2}}{x ^ {2} + y ^ {2}}, & (x, y) \neq (0, 0); \\ 0, & (x, y) = (0, 0). \end{array} \right.}
% p171-b05; source PDF page 171
\blockanchor{p171-b05}
则  $\frac{\partial^{2}f}{\partial x\partial y}(0,0)+\frac{\partial^{2}f}{\partial y\partial x}(0,0)=$  \_\_\_\_.\par
% p171-b06; source PDF page 171
\blockanchor{p171-b06}
\question{7}Laplace 算子  $\Delta f = \frac{\partial^{2}f}{\partial x^{2}} + \frac{\partial^{2}f}{\partial y^{2}}$  在极坐标变换  $\left\{\begin{aligned}x &= r\cos\theta \\ y &= r\sin\theta\end{aligned}\right.$$(0 \leq \theta < 2\pi)$  下的表达式为 \_\_\_\_.\par
% p171-b07; source PDF page 171
\blockanchor{p171-b07}
\question{8}$(\partial \mathbb{Q}^2)^\circ = \underline{\quad}$.\par
% p171-b08; source PDF page 171
\blockanchor{p171-b08}
\question{9}$\iint_{x^2 + y^2 \leq 4} (x^2 \sin y + x \ln (1 + y^2)) \, \mathrm{d}x \, \mathrm{d}y =$ \_\_\_\_.\par
% p171-b09; source PDF page 171
\blockanchor{p171-b09}
\question{10}设 C 为心脏线  $r = 1 + \cos \theta$ ，积分  $\oint_{C} (x + y^{2}) \, \mathrm{d}x + (y + x^{2}) \, \mathrm{d}y =$  \_\_\_\_.\par
% p171-b10; source PDF page 171
\par\Needspace{5\baselineskip}
\blockanchor{p171-b10}
\majorquestion{二}解答题. 每题 15 分, 共 90 分.\par
% p171-b11; source PDF page 171
\blockanchor{p171-b11}
\question{1}设 $f(x)$ 在 $[0, +\infty)$ 上一致连续, 证明: $f(x) = O(x) (x \to +\infty)$.\par
% p171-b12; source PDF page 171
\par\Needspace{6\baselineskip}
\blockanchor{p171-b12}
\question{2}设幂级数 $\sum_{n=0}^{\infty} a_n x^n$ 在 $(-1, 1)$ 上收敛, 且 $a_n \geq 0 (\forall n \geq 0)$. 又设\par
\nopagebreak[4]
% p171-b13; source PDF page 171
\blockanchor{p171-b13}
\displayformula{\lim _ {x \rightarrow 1 ^ {-}} \sum_ {n = 0} ^ {\infty} a _ {n} x ^ {n} = S \in \mathbb {R}.}
% p171-b14; source PDF page 171
\blockanchor{p171-b14}
证明: $\sum_{n=0}^{\infty} a_n$ 收敛于 $S$.\par
% p171-b15; source PDF page 171
\blockanchor{p171-b15}
\question{3}设  $x_{n}=1+\frac{1}{\sqrt{2}}+\cdots+\frac{1}{\sqrt{n}}-2\sqrt{n}(\forall n\geq1)$ ，证明： $\{x_{n}\}$  收敛.\par
% p171-b16; source PDF page 171
\blockanchor{p171-b16}
\question{4}设 $\Gamma(a) = \int_{0}^{+\infty} x^{a-1} e^{-x} \, \mathrm{d}x$, $a > 0$. 证明:\par
% p171-b17; source PDF page 171
\blockanchor{p171-b17}
(1) $\Gamma(a + 1) = a\Gamma(a), a > 0.$\par
% p171-b18; source PDF page 171
\blockanchor{p171-b18}
(2) $\ln \Gamma(a)$ 是凸函数.\par
% p171-b19; source PDF page 171
\par\Needspace{6\baselineskip}
\blockanchor{p171-b19}
\question{5}设 $E$ 是 $\mathbb{R}^n$ 中的开集, 函数 $f: E \to \mathbb{R}^m$ 与 $g: E \to \mathbb{R}^m$ 均可微. 设 $h(\pmb{x}) = \langle f(\pmb{x}), g(\pmb{x}) \rangle$ 为 $f(\pmb{x})$ 与 $g(\pmb{x})$ 的内积, 证明: $h(\pmb{x})$ 在 $E$ 上可微, 且对任意的 $\pmb{x} \in E$, 有\par
\nopagebreak[4]
% p171-b20; source PDF page 171
\blockanchor{p171-b20}
\displayformula{h ^ {\prime} (\boldsymbol {x}) = f (\boldsymbol {x}) ^ {T} g ^ {\prime} (\boldsymbol {x}) + g (\boldsymbol {x}) ^ {T} f ^ {\prime} (\boldsymbol {x}).}
% p171-b21; source PDF page 171
\par\Needspace{6\baselineskip}
\blockanchor{p171-b21}
\question{6}设 $\Omega \subset \mathbb{R}^2$ 为非空的连通开集, $f: \Omega \to \mathbb{R}$ 为连续函数. 且\par
\nopagebreak[4]
% p171-b22; source PDF page 171
\blockanchor{p171-b22}
\displayformula{f (\boldsymbol {x}) = \frac {1}{\pi r ^ {2}} \int_ {B _ {r} (\boldsymbol {x})} f, \forall \overline {{B _ {r} (\boldsymbol {x})}} \subset \Omega .}
% p171-b23; source PDF page 171
\blockanchor{p171-b23}
这里 $B_{r}(\pmb{x})$ 表示以 $\pmb{x}$ 为中心, 半径为 $r$ 的开圆盘. 证明: 若 $f$ 在 $\Omega$ 上有最大值, 则 $f$ 为常值函数.\par
